\documentclass[11pt,letterpaper]{article}
\usepackage[margin=0.7in,headheight=15pt]{geometry}
\usepackage[T1]{fontenc}\usepackage{lmodern,amsmath,amssymb,fancyhdr,parskip,hyperref}
\hypersetup{colorlinks=true,urlcolor=blue}
\pdfinfoomitdate=1\pdftrailerid{}\pdfsuppressptexinfo=15
\pagestyle{fancy}\fancyhf{}\lhead{Week 82 / Endless queues}\rhead{Facilitator draft}\cfoot{Bellingham Math Circle / Week 82 / facilitator / \thepage}
\setlength{\parskip}{7pt}\setlength{\parindent}{0pt}
\begin{document}
\section*{Mathematical overview}
An order-preserving matching is a bijection that preserves before/after relations. For the endless queue $1,2,3,\ldots$, prepending a star admits the matching star$\mapsto1$, $n\mapsto n+1$. Thus the order type $1+\omega$ is $\omega$. Appending a star after \emph{every} numbered customer produces a last element; the original queue has none. Hence $\omega+1$ is a different order type.

Two consecutive endless blocks have type $\omega+\omega$. The first customer in the second block has infinitely many predecessors and no immediate predecessor. An endless row of pairs has type $2\cdot\omega=\omega$; every customer has finitely many predecessors. They cannot be matched preserving order. All these queues have the same countable cardinal size, using explicit matches that need not preserve order.

Children should discover a structural obstruction or a rule covering every customer. Finite drawn matches are experiments; a uniform assignment and inverse establish an infinite matching. Grades 3--5 only, with facility interpreting a finite schematic for an endless rule. Ordinal addition here means concatenation. Surreal addition is a distinct commutative operation, so do not transfer these addition identities to Week 83's numbers.
\newpage
\section*{Preparation, launch and readiness}
Per pair: numbered customer cards, star, two colors of continuation strips, reusable queue boards and yarn or dry-erase arrows. A continuation symbol is an instruction, never another customer. A star after every number uses a \emph{separate panel} so children do not read it as after just the printed prefix. Colors and labels identify customers; they do not restrict which customers may be partners. Problem 4's small boxes are sketch space: arrange the larger cutouts on the table and sketch the result, rather than fitting the strips inside those boxes. Provide sufficient table space for two boards together; enlarging to two joined Letter sheets is possible. Physical fit untested.

Prerequisites: compare positions, follow a matching rule, understand ``each customer has exactly one partner.'' Reading is brief and can be read aloud. Reasoning load lies in extending a rule to an arbitrary named position, not counting a printed finite list. One anchored adult can coach pairs; at the upper triplet rotate constructor/checker/questioner.

Launch together: let children reorder a few finite cards, then show a finite matching with input, one middle pair and all completed pairs. Ask the partner to check both unused customers and crossing order. Introduce the continuing-strip convention \emph{before} asking about an endless queue.

First route after movement: ten minutes on finite Problem 1, 15 minutes prepending and matching, 15 minutes appending and explaining its obstruction. Problem 3's two blocks versus pairs is a second visit or encore; Problem 4 is open construction when the basic convention is controlled. Ordinal names may be introduced after the need arises. Do not require all stations or notation in one hour. Whole-group launch, then selected participating tables only.
\newpage
\section*{Facilitator explanations and hints}
Problem 1: list the B partners of R1, R2, R3. The six lists are\[ (1,2,3),(1,3,2),(2,1,3),(2,3,1),(3,1,2),(3,2,1). \] Only the first keeps every pair in order. Exhaustiveness follows from three choices, then two, then one.

Problem 2, prepending: the assignment is onto because customer 1 receives the star and customer $m\ge2$ receives original $m-1$. It is one-to-one and preserves order. Matching only six physical cards leaves an apparent unmatched card at the right edge; the repeating rule, not this artificial drawing boundary, resolves the infinite case.

Appending: suppose there were an order-preserving bijection to the ordinary endless queue. The image of the appended last customer would have to be last, but every ordinary numbered customer has a successor. Contradiction. The issue is not inability to physically walk past infinitely many customers.

Ignoring order, star$\mapsto1$, $n\mapsto n+1$ still gives a size match for the appended-star queue. It moves the last customer to the front, showing why the preservation condition matters. It does not make the queues have equal order type.

Problem 3: D has all R customers followed by all B customers; E interleaves R and B. Matching each named R$_n$ and B$_n$ in D to its namesake in E gives an unrestricted size match. Their positions in E are $2n-1$ and $2n$. This interleaves blocks and violates their original order. No order-preserving match can send B$_1$ to a customer with only finitely many predecessors, since B$_1$ has all R customers before it.

Problem 4 examples: a row alone versus a star followed by the row admits the order-preserving shift of Problem 2. A row alone versus a row followed by a star admits a size match but no order match because of the last customer. Each construction uses one row and at most two stars.

Hints: point to a last customer if one exists; ask about the first customer of block B; ask how many customers are before a \emph{named} position. Avoid directing every match in advance. The representation is conceptual; endless strips need not physically contain every customer.
\newpage
\section*{Encore, sources and limits}
Encore: attach two, then three finite customers before or after the endless queue. Compare removing the first customer with removing an appended last customer. Later arrange endlessly many endless blocks as rows with explicit row order; this opens a doorway toward $\omega^2$ but is not required for the present proof.

Sources: Joel David Hamkins, ``Counting to Infinity and Beyond'' (September 2, 2021), \url{https://jdh.hamkins.org/counting-to-infinity-poster/}, for ordinal visualization and continuation to $\omega^2$. Our matching boards and explanatory sequence are independently authored; no poster or video frame is republished. Basic order-type arguments are given fully in this guide.

Pedagogy: \emph{Math Circle by the Bay}, preface p. viii, connects accessible topics to deep mathematics; pp. ix--x discuss materials and varying depth. The grade gate here is our unpiloted design judgment. Existing Week 68 endless graphs explores ends, not these queue order types.

Unpiloted. Print US Letter at Actual Size. Physical board fit, yarn handling, table-space requirement, timing and classroom procedure have not been rehearsed. Record whether children interpreted continuation as a rule, whether they distinguished size/order, and any useful invented queues. A decorative infinity symbol alone is not evidence that the child controlled the investigation.
\end{document}
